%-----------------------------------------------------------------------------------------------------------------------------------------------%
%	The MIT License (MIT)
%
%	Copyright (c) 2021 Jitin Nair
%
%	Permission is hereby granted, free of charge, to any person obtaining a copy
%	of this software and associated documentation files (the "Software"), to deal
%	in the Software without restriction, including without limitation the rights
%	to use, copy, modify, merge, publish, distribute, sublicense, and/or sell
%	copies of the Software, and to permit persons to whom the Software is
%	furnished to do so, subject to the following conditions:
%	
%	THE SOFTWARE IS PROVIDED "AS IS", WITHOUT WARRANTY OF ANY KIND, EXPRESS OR
%	IMPLIED, INCLUDING BUT NOT LIMITED TO THE WARRANTIES OF MERCHANTABILITY,
%	FITNESS FOR A PARTICULAR PURPOSE AND NONINFRINGEMENT. IN NO EVENT SHALL THE
%	AUTHORS OR COPYRIGHT HOLDERS BE LIABLE FOR ANY CLAIM, DAMAGES OR OTHER
%	LIABILITY, WHETHER IN AN ACTION OF CONTRACT, TORT OR OTHERWISE, ARISING FROM,
%	OUT OF OR IN CONNECTION WITH THE SOFTWARE OR THE USE OR OTHER DEALINGS IN
%	THE SOFTWARE.
%	
%
%-----------------------------------------------------------------------------------------------------------------------------------------------%

%----------------------------------------------------------------------------------------
%	DOCUMENT DEFINITION
%----------------------------------------------------------------------------------------

% article class because we want to fully customize the page and not use a cv template
\documentclass[a4paper,11pt]{article}

%----------------------------------------------------------------------------------------
%	FONT
%----------------------------------------------------------------------------------------

% % fontspec allows you to use TTF/OTF fonts directly
% \usepackage{fontspec}
% \defaultfontfeatures{Ligatures=TeX}

% % modified for ShareLaTeX use
% \setmainfont[
% SmallCapsFont = Fontin-SmallCaps.otf,
% BoldFont = Fontin-Bold.otf,
% ItalicFont = Fontin-Italic.otf
% ]
% {Fontin.otf}

%----------------------------------------------------------------------------------------
%	PACKAGES
%----------------------------------------------------------------------------------------
\usepackage{url}
\usepackage{parskip} 	

%other packages for formatting
\RequirePackage{color}
\RequirePackage{graphicx}
\usepackage[usenames,dvipsnames]{xcolor}
\usepackage[scale=0.94]{geometry}

%tabularx environment
\usepackage{tabularx}

%for lists within experience section
\usepackage{enumitem}

% centered version of 'X' col. type
\newcolumntype{C}{>{\centering\arraybackslash}X} 

%to prevent spillover of tabular into next pages
\usepackage{supertabular}
\usepackage{tabularx}
\newlength{\fullcollw}
\setlength{\fullcollw}{0.47\textwidth}

%custom \section
\usepackage{titlesec}				
\usepackage{multicol}
\usepackage{multirow}

%CV Sections inspired by: 
%http://stefano.italians.nl/archives/26
\titleformat{\section}{\Large\scshape\raggedright}{}{0em}{}[\titlerule]
\titlespacing{\section}{0pt}{11pt}{8pt}

%for publications
\usepackage[style=authoryear,sorting=ynt, maxbibnames=2]{biblatex}

%Setup hyperref package, and colours for links
\usepackage[unicode, draft=false]{hyperref}
\definecolor{linkcolour}{rgb}{0,0.2,0.6}
\hypersetup{colorlinks,breaklinks,urlcolor=linkcolour,linkcolor=linkcolour}
\addbibresource{citations.bib}
\setlength\bibitemsep{0.35em}

% Compact global spacing so long CV content fits in two pages.
\setlength{\parskip}{0.45em}
\setlength{\tabcolsep}{6pt}
\renewcommand{\arraystretch}{1.08}

\makeatletter
\renewcommand{\today}{\number\day\ de\ \ifcase\month\or enero\or febrero\or marzo\or abril\or mayo\or junio\or julio\or agosto\or septiembre\or octubre\or noviembre\or diciembre\fi\ de\ \number\year}
\makeatother

%for social icons
\usepackage{fontawesome5}

%debug page outer frames
%\usepackage{showframe}


% job listing environments
\newenvironment{jobshort}[2]
    {
    \begin{tabularx}{\linewidth}{@{}l X r@{}}
    \textbf{#1} & \hfill &  #2 \\[3.75pt]
    \end{tabularx}
    }
    {
    }

\newenvironment{joblong}[2]
    {
    \begin{tabularx}{\linewidth}{@{}l X r@{}}
    \textbf{#1} & \hfill &  #2 \\[3.75pt]
    \end{tabularx}
    \begin{minipage}[t]{\linewidth}
    \begin{itemize}[nosep,after=\strut, leftmargin=1em, itemsep=3pt,label=--]
    }
    {
    \end{itemize}
    \end{minipage}    
    }



%----------------------------------------------------------------------------------------
%	BEGIN DOCUMENT
%----------------------------------------------------------------------------------------
\begin{document}

% non-numbered pages
\pagestyle{empty} 

%----------------------------------------------------------------------------------------
%	TITLE
%----------------------------------------------------------------------------------------

% \begin{tabularx}{\linewidth}{ @{}X X@{} }
% \huge{Your Name}\vspace{2pt} & \hfill \emoji{incoming-envelope} email@email.com \\
% \raisebox{-0.05\height}\faGithub\ username \ | \
% \raisebox{-0.00\height}\faLinkedin\ username \ | \ \raisebox{-0.05\height}\faGlobe \ mysite.com  & \hfill \emoji{calling} number
% \end{tabularx}

\begin{tabularx}{\linewidth}{@{} C @{}}
\Huge{Matías Ignacio Serrano Acevedo} \\[7.5pt]

%\href{https://linkedin.com/in/username}{\raisebox{-0.05\height}\faLinkedin\ username} \ $|$ \ 
\href{https://misaa.cc}{\faGlobe \ misaa.cc} \ $|$ \
\href{https://github.com/misaaaaaa}{\faGithub\ misaaaaaa} \ $|$ \ 
\href{mailto:matias.serranoacevedo@gmail.com}{\faEnvelope \ matias.serranoacevedo@gmail.com} \ $|$ \
{Santiago, Chile} \\
\end{tabularx}



%----------------------------------------------------------------------------------------
% EXPERIENCE SECTIONS
%----------------------------------------------------------------------------------------

%----------------------------------------------------------------------------------------
%	EDUCATION
%----------------------------------------------------------------------------------------
\section{Formación}
\begin{tabularx}{\linewidth}{@{}l X@{}}	
2026 - \_\_\_\_ & Doctorado en Artes y Humanidades - \textbf{Universidad de Santiago de Chile} \hfill \normalsize  \\

2021 - 2024 & Magíster en Artes Mediales - \textbf{Universidad de Chile} \hfill \normalsize  \\

2011 - 2017 & Licenciatura en Artes con mención en Sonido - \textbf{Universidad de Chile} \hfill  \\ 

%2022 & Class 12th Some Board \hfill  (Grades) \\

%2021 & Class 10th Some Board \hfill  (Grades) \\
\end{tabularx}

%Interests/ Keywords/ Summary
%\section{Summary}
%This CV can also be automatically complied and published using GitHub Actions. For details, \href{https://github.com/jitinnair1/autoCV}{click here}.

%Experience
\section{Proyectos de Artes Visuales}
\begin{tabularx}{\linewidth}{@{}l X@{}}	

2026 & \textit{Llluvia Metropolitana}. Instalación sonora junto a Rainer Krause @  \textbf{Galería Metroplitana}. Santiago, Chile. \hfill \normalsize \href{https://misaa.cc/projects/llluviametropolitana.html.html}{info} \\

2025 & \textit{oí 1, oí 2 \& oí 3}. Instalación sonora en exposición colectiva con Almendra Díaz @  \textbf{Galería de Artes Visuales UNIACC}. Santiago, Chile. \hfill \normalsize \href{https://misaa.cc/projects/cuandotodaslaspalabras.html}{info} \\

2025 & \textit{Hablar con extraños}. Instalación sonora @  \textbf{Galería Micromedios}. Santiago, Chile. \hfill \normalsize \href{https://translab.uchilefau.cl/hablar-con-extranos/}{info} \\

2025 & \textit{Un trazo en el oído}. Instalación sonora en exposición colectiva "Tiempo de Decaimiento Temprano" @  \textbf{Museo de Arte Contemporáneo}. Santiago, Chile. \hfill \normalsize \href{https://misaa.cc/projects/untrazoeneloido.html/}{info} \\
2025 & \textit{Transistores en tránsito}. Instalación sonora en exposición colectiva "Tiempo de Decaimiento Temprano" @  \textbf{Museo de Arte Contemporáneo}. Santiago, Chile. \hfill \normalsize \href{https://misaa.cc/projects/transistoresentransito.html}{info} \\

2024 & \textit{Amplificar la duda}. Instalación sonora en encuentro "VOLCÁN" @  \textbf{MIM}. Santiago, Chile. \hfill \normalsize \href{https://mim.cl/noticia/matias-serrano-escuchar-necesariamente-tiene-que-ver-con-el-presente-y-el-aqui-y-el-ahora/}{interview} \\

2024 & \textit{Museo en Estéreo}. Instalación sonora creada después de un taller con vecinos del edificio de edificios del museo @  \textbf{Museo de Arte Contemporáneo}. Santiago, Chile. \hfill \normalsize \href{https://misaa.cc/projects/museoenestereo.html}{info} \\

2025 & \textit{Lenguaje y comunicación}. Instalación sonora en exposición colectiva "Balmaceda visual: Bordes Fluidos" @  \textbf{Museo de Arte Contemporáneo}. Santiago, Chile. \hfill \normalsize \href{https://misaa.cc/projects/lenguajeycomunicacion.html}{info} \\

2023 & \textit{La memoria contándole a la luz el sonido de la lluvia}.  Instalación sonora en exposición colectiva "XVI Premio MAVI-UC". \textsc{audience award} @  \textbf{Museo de Artes Visuales}. Santiago, Chile. \hfill \normalsize \href{https://misaa.cc/projects/lamemoria.html}{info} \\

2019 & \textit{Dispositivas de encarnación}. Co-created as Colectiva 22bits. Instalación sonora en exposición colectiva "Marea: Arte y espacio marítimo" @ \textbf{CENTEX, Ministry of Culture}. Valparaíso, Chile. \hfill \normalsize \href{https://centex.cultura.gob.cl/wp-content/uploads/2019/07/catalogo_marea_digital.-1.pdf}{exhibition catalog} \\

2018 & \textit{Espectral: Escucha de un paisaje intervenido}. Co-created as Colectiva 22bits. Instalación sonora @  \textbf{Casa 916}. Concepción, Chile. \hfill \normalsize \href{https://misaa.cc/projects/espectral.html}{info} \\

2017 & \textit{Ensayo de horizontalidad}. Co-created as Colectiva 22bits. Instalación sonora en exposición colectiva "6° Concurso de Arte y Tecnologías Digitales en Homenaje a Matilde Pérez". \textsc{3rd place, students tier} @  \textbf{Museo de Arte Contemporáneo}. Santiago, Chile. \hfill \normalsize \href{https://misaa.cc/projects/edh.html}{info} \\

\end{tabularx}

\section{Participación en proyectos interdisciplinarios y colectivos}
\begin{tabularx}{\linewidth}{@{}l X@{}}

2023 - \_\_\_\_ & \textit{Núcleo de Artes Sonoras}. Investigación y promoción de las artes sonoras dentro de las Artes Visuales. Departamento de Artes Visuales, U. de Chile. \\

2025 - 2026 & \textit{Aproximaciones subacuáticas}. Proyecto de creación e investigación interdisciplinar (arte sonoro y nado). Proyecto Fondart dirigido por Fernanda Fábrega. \\

2024 & \textit{Árboles ciudadanos}. Proyecto de creación e investigación interdisciplinar (electrónica, sonido, teatro y espacio público) en torno a los árboles en entornos urbanos. Producido por compañía Organismo Teatro. \\

2020 & \textit{Ruido propio}. Núcleo de investigación de memoria e identidad sonora dirigido por Dra. Carla Badani. Departamento de Sonido, Universidad de Chile. \\

2019 - 2024 & \textit{Archivo Veintidós}. Sello de Música Experimental y Arte Sonoro del sur global co-creado y co-dirigido con Bárbara Molina (Colectiva 22bits). \\

2018 - 2020 & \textit{Casa Oram}. Espacio de talleres y actividades en torno a la tecnología y la escucha ubicado en barrio Matta Sur, coordinado en conjunto a Bárbara Molina (Colectiva 22bits). \\

2016 & \textit{Tutupá}. Proyecto de investigación y creación tecnológica. Convocatoria Becas Santiago Maker Space. \\

\end{tabularx}

\newpage
\small
\setlength{\parskip}{0.05em}
\renewcommand{\arraystretch}{0.9}
\titlespacing{\section}{0pt}{6pt}{3pt}

%\begin{jobshort}{Designation}{Jan 2021 - present}
%long long line of blah blah that will wrap when the table fills the column width long long line of blah blah that will wrap when the table fills the column width long long line of blah blah that will wrap when the table fills the column width long long line of blah blah that will wrap when the table fills the column width
%\end{jobshort}


%\begin{joblong}{Designation}{Mar 2019 - Jan 2021}
%\item long long line of blah blah that will wrap when the table fills the column width
%\item again, long long line of blah blah that will wrap when the table fills the column width but this time even more long long line of blah blah. again, long long line of blah blah that will wrap when the table fills the column width but this time even more long long line of blah blah
%\end{joblong}

\section{Residencias}
\begin{tabularx}{\linewidth}{@{}l X@{}}	

Enero, 2025 & \textit{Radicante}. Residencia en un barco. Organizada por \textsc{Liquen LAB} @  \textbf{Estrecho de Magallanes}. Punta Arenas, Chile. \hfill \normalsize \href{https://www.instagram.com/p/DFvAEfCok0J/}{testimonial} \\

Septiembre, 2024 & \textit{Proyecto ERROR}. Residencia de arte contemporáneo. @  \textbf{Ciudad de México}. CDMX, México. \hfill \normalsize \href{https://www.instagram.com/p/C_3IjYUOlui}{welcoming} \\

Marzo, 2023 & \textit{Utopías y distopías en la modernidad}. Residencia con la artista chilena Ingrid Wildi para crear una banda sonora para una película inconclusa @  \textbf{Isla King George}. Antártica. \hfill \normalsize \href{https://misaa.cc/projects/utopiasydistopias.html}{info} \\

Noviembre, 2019 & \textit{Encuentro Tsonami}. Colectiva 22bits. Prácticas sonoras en contextos de crisis @  \textbf{Tsonami Sound Art Festival}. Valparaíso, Chile. \hfill \normalsize \href{https://misaa.cc/projects/tsonami2019.html}{info} \\

Agosto, 2019 & \textit{La Jacquer EsCool}. Colectiva 22bits. Residencia de intercambio de conocimientos @  \textbf{Platohedro}. Medellín, Colombia. \hfill \normalsize \href{https://platoteca.platohedro.org/?p=9508}{entrevista} \\

Septiembre, 2018 & \textit{Residencia de creación}. Colectiva 22bits. Investigación y producción de "Dispositivas de Encarnación" @  \textbf{Sala B.A.S.E., Tsonami}. Valparaíso, Chile. \hfill \normalsize \href{https://misaa.cc/projects/dispositivas.html}{info} \\

\end{tabularx}
  
%Projects
\section{Experiencia docente}

\begin{tabularx}{\linewidth}{ @{}l r@{} }
\textbf{Universidad de Chile} & \hfill  \\[3.75pt]
\multicolumn{2}{@{}X@{}}{\textit{Académico adjunto, Departamento de Artes Visuales} 


Medios II (Introducción a las Artes Mediales), desde 2022

Taller complementario de Arte y Electrónica, desde 2023

Taller de Electrónica y Programación, desde 2023}  \\
\end{tabularx}

\begin{tabularx}{\linewidth}{ @{}l r@{} }
\textbf{Universidad Diego Portales (UDP)} & \hfill  \\[3.75pt]
\multicolumn{2}{@{}X@{}}{\textit{Profesor externo} 

Taller optativo de Arte y Electrónica (Escuela de Artes Visuales), desde 2024

Diseño de máquinas electrónicas (Escuela de Diseño, junto a Aarón Montoya), desde 2025

Diseño de máquinas computacionales (Escuela de Diseño, junto a Aarón Montoya), desde 2025

Pensamiento Computacional (Escuela de Diseño), desde 2026}  
\\
\end{tabularx}

\begin{tabularx}{\linewidth}{ @{}l r@{} }
\textbf{Universidad de Artes, Ciencias y Comunicaciones (UNIACC)} & \hfill  \\[3.75pt]
\multicolumn{2}{@{}X@{}}{\textit{Profesor externo} 

Tecnología Aplicada I: Electrónica, desde 2023

Tecnología Aplicada III: Objeto, desde 2024
}  
\\
\end{tabularx}

%----------------------------------------------------------------------------------------
%	PUBLICATIONS
%----------------------------------------------------------------------------------------
\section{Publicaciones}
\begin{refsection}[citations.bib]
\nocite{*}
\printbibliography[heading=none]
\end{refsection}

%----------------------------------------------------------------------------------------
%	SKILLS
%----------------------------------------------------------------------------------------
\section{Becas y premios}

\begin{tabularx}{\linewidth}{@{}l X@{}}
2026 &  \normalsize{\textsc{Beca Completa}. Beca interna de arancel y mantención para formar parte del programa de Doctorado en Artes y Humanidades de la Universidad de Santiago de Chile.}\\
\end{tabularx}

\begin{tabularx}{\linewidth}{@{}l X@{}}
2023 &  \normalsize{\textsc{Premio del público}. Obra: 'La memoria contándole a la luz el sonido de la lluvia'. XVI Premio de Arte Joven MAVI-UC.}\\
\end{tabularx}

\begin{tabularx}{\linewidth}{@{}l X@{}}
2019 &  \normalsize{\textsc{Beca Chile Crea}. Proyecto: “Tutorías de Electrónica para Nuevos Medios y Arte Sonoro,” guiado por el profesor José Luis Cárdenas. Ministerio de las Artes, las Culturas y el Patrimonio}\\
\end{tabularx}

\begin{tabularx}{\linewidth}{@{}l X@{}}
2018 &  \normalsize{\textsc{Beca de creación y producción}. Proyecto: 'Espectral: Escucha de un paisaje intervenido'. Ministerio de las Artes, las Culturas y el Patrimonio}\\
\end{tabularx}

\begin{tabularx}{\linewidth}{@{}l X@{}}
2017 &  \normalsize{\textsc{3er lugar, categoría estudiantes}. Obra: 'Ensayo de Horizontalidad'. 6to Concurso de Artes y Tecnologías Digitales – Homenaje a Matilde Pérez}\\
\end{tabularx}

\vfill
\center{\footnotesize Última actualización: \today}

\end{document}
